\documentclass[11pt,a4paper]{article}

\usepackage[T2A]{fontenc}
\usepackage[utf8]{inputenc}
\usepackage[russian]{babel}
\usepackage[a4paper,margin=20mm]{geometry}
\usepackage{microtype}

\usepackage{amsmath,amssymb,bm,mathtools}
\usepackage{booktabs}
\usepackage{array}
\usepackage{xcolor}
\usepackage{hyperref}
\hypersetup{
    colorlinks=true,
    linkcolor=blue!60!black,
    urlcolor=blue!60!black,
    citecolor=blue!60!black
}

\emergencystretch=1.5em

\title{\textbf{Квазиодномерная нестационарная модель нагрева солнечной короны и ускорения солнечного ветра альфвеновской турбулентностью (AWSoM-1D):\\ обзор литературы, математический анализ и численные методы}}
\author{}
\date{}

\begin{document}

\maketitle

\begin{abstract}
В работе представлена замкнутая математическая формулировка, параметрический анализ и вычислительная методология решения квазиодномерной нестационарной системы уравнений магнитной гидродинамики (МГД) с самосогласованным переносом, линейным отражением и каскадной диссипацией альфвеновских волн вдоль расширяющейся магнитной силовой трубки (одномерная редукция модели AWSoM). Детально рассмотрена физическая модель переходного слоя между хромосферой и короной, включающая анизотропную электронную теплопроводность Спитцера--Харма, радиационные потери оптически тонкой плазмы и аналитическую регуляризацию теплового фронта Линкера--Лионелло--Соколова. Выполнен строгий характеристический анализ системы: показано, что присутствие модуля пространственного градиента альфвеновской скорости в коэффициенте отражения $\mathcal{R}$ превращает правую часть волновых уравнений на градиентной ветви в неконсервативный дифференциальный оператор первого порядка, модифицирующий матрицу Якоби и характеристические скорости. Систематизированы фундаментальные вычислительные трудности задачи (многомасштабная характеристическая и источниковая жёсткость, неконсервативные произведения, корневая особенность $\sqrt{w_+ w_-}$ при вырождении встречной волны, негладкость операторов $\min/\max$ и тепловой пограничный слой) и построена безусловно положительная схема метода установления, сочетающая метод конечного объёма MUSCL-HLL/Rusanov, локальное полунеявное расщепление в переменных амплитуд и неявную трёхдиагональную прогонку Томаса.
\end{abstract}

\tableofcontents
\newpage

% ==============================================================================
\section{Введение и обзор литературы}
% ==============================================================================

Проблема механизмов нагрева солнечной короны до температур $T \sim (1\text{--}3)\times 10^6$~К и последующего разгона сверхзвукового и суперальфвеновского солнечного ветра остаётся одной из центральных задач физики космической плазмы со времени пионерской работы Ю.~Паркера \cite{Parker1958}. Классическая изотермическая и политропная модели Паркера продемонстрировали неизбежность трансзвукового расширения горячей короны в вакуум, однако оставили открытым вопрос о природе распределённого источника энергии и импульса, способного преодолеть солнечную гравитацию и сформировать наблюдаемые потоки быстрого ($u_\infty \approx 700\text{--}800$~км/с) и медленного ($u_\infty \approx 300\text{--}450$~км/с) солнечного ветра.

Уже в конце 1960-х и 1970-х годах наблюдения межпланетных мерцаний и прямые измерения на космических аппаратах \emph{Mariner~2} и \emph{Helios} обнаружили присутствие в солнечном ветре мощных флуктуаций магнитного поля и скорости, обладающих высокой корреляцией $\delta\bm{u} \approx \mp \delta\bm{B}/\sqrt{\mu_0\rho}$, характерной для альфвеновских волн, распространяющихся от Солнца \cite{BelcherDavis1971}. Это привело к формированию парадигмы \textbf{волнового и турбулентного ускорения и нагрева солнечного ветра}, развитие которой прошло несколько фундаментальных этапов.

\subsection{Приближение ВКБ и пондеромоторное волновое давление}

На первом этапе в работах Алазраки и Кутурье \cite{Alazraki1971}, Белчера \cite{Belcher1971}, Холлвега \cite{Hollweg1973,Hollweg1978} и Жака \cite{Jacques1977} была построена теория распространения коротковолновых альфвеновских волн в неоднородном движущемся потоке плазмы в приближении геометрической оптики (ВКБ, $\lambda \ll L_A = |\nabla\ln V_A|^{-1}$). Было строго показано, что усреднённое по периоду воздействие альфвеновских флуктуаций на основное МГД-течение сводится к появлению в уравнении импульса градиента изотропного волнового давления $P_A = (w_+ + w_-)/2$, а в уравнении переноса плотности волновой энергии $w_+$ --- характерного конвективно-рабочего члена с коэффициентом $3/2$:
\[
\nabla\cdot\left[(\bm{u} + \bm{V}_A)w_+\right] + \frac{w_+}{2}\nabla\cdot\bm{u} = 0,
\]
который в отсутствие диссипации сохраняет волновое действие (адиабатический инвариант Блоха--Бретертона--Девара). В рамках чистого ВКБ-приближения волны передают фоновому потоку только механическую работу (разгоняя ветер за счёт силы $-\nabla P_A$), но не выделяют тепла и не порождают встречных волн $w_-$.

\subsection{Не-ВКБ отражение волн и запуск МГД-каскада}

Фундаментальное ограничение приближения ВКБ состоит в том, что в несжимаемой идеальной МГД нелинейное взаимодействие и каскадный перенос энергии в малые поперечные масштабы ($\sim k_\perp$) возможны \emph{только при одновременном присутствии встречных альфвеновских волновых пакетов}, описываемых переменными Эльзассера $\bm{z}_\pm = \delta\bm{u} \mp \delta\bm{B}/\sqrt{\mu_0\rho}$ \cite{Kraichnan1965,Dobrowolny1980}:
\[
\frac{\partial\bm{z}_\pm}{\partial t} + (\bm{u} \pm \bm{V}_A)\cdot\nabla\bm{z}_\pm + (\bm{z}_\mp\cdot\nabla)\bm{z}_\pm = \dots
\]
Если на нижней границе Солнца генерируются только уходящие волны ($\bm{z}_+ \neq \bm{0}$, $\bm{z}_- = \bm{0}$), то нелинейный член $(\bm{z}_-\cdot\nabla)\bm{z}_+$ тождественно равен нулю, и турбулентный нагрев не возникает.

Механизм генерации затравочной встречной волны $\bm{z}_-$ в открытых магнитных трубках был установлен в классической работе Хайнеманна и Олберта \cite{HeinemannOlbert1980}, а также в работах Велли \cite{Velli1993} и Маттеуса и др. \cite{Matthaeus1999}: в стратифицированной короне с резким радиальным градиентом плотности и магнитного поля длина волны низкочастотных альфвеновских колебаний оказывается сопоставимой с масштабом неоднородности альфвеновской скорости $V_A(r)$. Это вызывает частичное \textbf{линейное (не-ВКБ) отражение} выходящих волн на градиентах фона со скоростью, пропорциональной $|(\bm{V}_A\cdot\nabla)\ln V_A| = |\partial V_A/\partial r|$. Отражённые волны $w_-$ вступают в нелинейное взаимодействие с основной популяцией $w_+$, запуская анизотропный перпендикулярный каскад энергии до ионных гирорадиусов, где энергия диссипирует в тепло плазмы \cite{GoldreichSridhar1995,Dmitruk2002}.

\subsection{Одномерные модели вдоль магнитной трубки и семейство AWSoM}

Начиная с середины 2000-х годов, этот физический механизм (распространение $w_+$ $\to$ линейное отражение $\mathcal{R}$ $\to$ генерация $w_-$ $\to$ нелинейный каскад $\Gamma_\pm$ $\to$ нагрев $Q_w$ и ускорение $-\nabla P_A$) лёг в основу целого класса самосогласованных одномерных (Flux-Tube) моделей солнечного ветра:
\begin{itemize}
\item В работах Сузуки и Инуцуки \cite{SuzukiInutsuka2005} и Кранмера, ван Баллегуйена и Эдгара \cite{Cranmer2007} были построены эталонные одномерные модели вдоль сверхрадиально расширяющихся магнитных трубок $A(r)$ от фотосферы/хромосферы до гелиосферы. Было показано, что задание единственного спектра волн у основания трубки позволяет одновременно воспроизвести как быстрый ветер в корональных дырах (большое сверхрадиальное расширение $f_{\max}$), так и медленный ветер на границах стримеров.
\item В работах Вердини и Велли \cite{VerdiniVelli2007,Verdini2010} и Чандрана и др. \cite{Chandran2009,Chandran2011} одномерные уравнения переноса волн были замкнуты в форме двух уравнений для плотностей энергии $w_+$ и $w_-$ с феноменологическими частотами каскадной диссипации Кармана--Хоуарта $\Gamma_\pm = \frac{2}{L_\perp}\sqrt{\frac{w_\mp}{\rho}}$, где поперечная корреляционная длина масштабируется с сохранением магнитного потока как $L_\perp \propto B^{-1/2} \propto A^{1/2}$ \cite{Hollweg1986}.
\item Наконец, в работах группы Мичиганского университета --- Соколова и др. \cite{Sokolov2013}, Оран и др. \cite{Oran2013} и ван дер Холста и др. \cite{vanderHolst2014} --- этот формализм был обобщён до глобальной трёхмерной модели солнечной короны и гелиосферы \textbf{AWSoM} (Alfv\'en Wave Solar Model). В работе \cite{vanderHolst2014} было введено рассматриваемое в настоящем документе самосогласованное замыкание для скорости отражения $\mathcal{R} = \min\{\mathcal{R}_{\rm imb}, \max(\Gamma_+,\Gamma_-)\}\mathcal{C}(w_+, w_-)$, которое учитывает подавление отражения при переходе от сильно несбалансированной турбулентности ($w_+ \gg w_-$) к сбалансированному состоянию ($w_+ \approx w_-$).
\item В последующих работах Соколова и др. \cite{Sokolov2016,Sokolov2021} (модель \textbf{AWSoM-R} / Threaded-Field-Line Model) уравнения AWSoM в нижней короне ($r \in [1, 1.15] R_\odot$) были редуцированы именно к квазиодномерной системе вдоль пучка магнитных силовых линий для сверхбыстрого расчёта переходного слоя и сшивки с трёхмерной МГД-областью.
\end{itemize}

Несмотря на широкое применение трёхмерного кода AWSoM (в рамках программного комплекса SWMF / BATS-R-US), в литературе до сих пор отсутствовало замкнутое математическое исследование самой квазиодномерной дифференциальной системы с замыканием ван дер Холста \cite{vanderHolst2014} во всей области от хромосферы до сверхзвукового ветра ($1 \le r/R_\odot \le 30$). Между тем эта система обладает нетривиальными математическими свойствами (скрытая квазилинейность источникового члена отражения по градиенту плотности, кусочно-гладкое переключение матрицы Якоби и экстремальная разномасштабность времён релаксации), детальному анализу и численному преодолению которых посвящена настоящая работа.

% ==============================================================================
\section{Исходная трёхмерная система и квазиодномерная редукция}
% ==============================================================================

\subsection{Трёхмерные уравнения модели AWSoM}

В трёхмерной постановке \cite{Sokolov2013,vanderHolst2014} состояние замагниченной плазмы и альфвеновской турбулентности в момент времени $t$ в точке $\bm{x}$ ($r = |\bm{x}|$) описывается плотностью $\rho$, массовой скоростью $\bm{u}$, изотропным тепловым давлением $p$, магнитной индукцией $\bm{B}$, а также плотностями энергии альфвеновских волн $w_+$ и $w_-$, распространяющихся вдоль и против $\bm{B}$.

Полная плотность энергии плазмы и магнитного поля (без энергии волн) $E_{\rm mhd}$, вектор скорости Альфвена $\bm{V}_A$, волновое давление $P_A$ и гравитационное ускорение $\bm{g}$ задаются соотношениями:
\begin{equation}
E_{\rm mhd} = \frac{\rho|\bm{u}|^2}{2} + \frac{p}{\gamma-1} + \frac{|\bm{B}|^2}{2\mu_0}, \qquad
\bm{V}_A = \frac{\bm{B}}{\sqrt{\mu_0\rho}}, \qquad
P_A = \frac{w_+ + w_-}{2}, \qquad
\bm{g} = -\frac{GM_\odot}{r^3}\bm{x},
\label{eq:defs_3d}
\end{equation}
где $\gamma = 5/3$. Система уравнений сохранения массы, импульса, магнитной индукции, энергии плазмы и переноса энергии волн имеет вид:
\begin{align}
\frac{\partial\rho}{\partial t} + \nabla\cdot(\rho\bm{u}) &= 0,
\label{eq:3d_mass}\\[4pt]
\frac{\partial(\rho\bm{u})}{\partial t} + \nabla\cdot\left[\rho\bm{u}\otimes\bm{u} + \left(p + P_A + \frac{|\bm{B}|^2}{2\mu_0}\right)\mathbf{I} - \frac{\bm{B}\otimes\bm{B}}{\mu_0}\right] &= \rho\bm{g},
\label{eq:3d_momentum}\\[4pt]
\frac{\partial\bm{B}}{\partial t} + \nabla\cdot(\bm{B}\otimes\bm{u} - \bm{u}\otimes\bm{B}) &= \bm{0}, \qquad \nabla\cdot\bm{B} = 0,
\label{eq:3d_induction}\\[4pt]
\frac{\partial E_{\rm mhd}}{\partial t} + \nabla\cdot\left[\left(E_{\rm mhd} + p + \frac{|\bm{B}|^2}{2\mu_0}\right)\bm{u} - \frac{(\bm{u}\cdot\bm{B})\bm{B}}{\mu_0} + \bm{q}\right] &= \rho\bm{u}\cdot\bm{g} - \bm{u}\cdot\nabla P_A + Q_w - Q_{\rm rad} + Q_{\rm bg},
\label{eq:3d_energy}\\[4pt]
\frac{\partial w_\pm}{\partial t} + \nabla\cdot\left[(\bm{u} \pm \bm{V}_A)w_\pm\right] + \frac{w_\pm}{2}\nabla\cdot\bm{u} &= \mp\mathcal{R}\sqrt{w_+ w_-} - \Gamma_\pm w_\pm.
\label{eq:3d_waves}
\end{align}
Здесь $\bm{q}$ --- поток электронной теплопроводности, $Q_{\rm rad}$ --- радиационные потери, $Q_{\rm bg}$ --- дополнительный фоновый нагрев (далее полагаемый равным нулю). Скорости каскадной диссипации $\Gamma_\pm$, мощность турбулентного нагрева $Q_w$ и коэффициент линейного отражения $\mathcal{R}$ замыкаются согласно \cite{vanderHolst2014}:
\begin{equation}
\Gamma_\pm = \frac{2}{L_\perp}\sqrt{\frac{w_\mp}{\rho}}, \qquad
L_\perp\sqrt{|\bm{B}|} = C_L = \mathrm{const} > 0, \qquad
Q_w = \Gamma_+ w_+ + \Gamma_- w_-,
\label{eq:3d_dissipation}
\end{equation}
\begin{equation}
\mathcal{R} = \min\left\{\mathcal{R}_{\rm imb},\, \max(\Gamma_+,\Gamma_-)\right\}\mathcal{C}(w_+, w_-), \qquad
\mathcal{R}_{\rm imb} = \sqrt{\left[(\bm{V}_A\cdot\nabla)\ln|\bm{V}_A|\right]^2 + \left[\bm{b}\cdot(\nabla\times\bm{u})\right]^2},
\label{eq:3d_reflection}
\end{equation}
где $\bm{b} = \bm{B}/|\bm{B}|$, а функция дисбаланса энергий встречных волн $\mathcal{C}(w_+, w_-)$ равна:
\begin{equation}
\mathcal{C}(w_+, w_-) =
\begin{cases}
1 - 2\sqrt{w_-/w_+}, & w_+ \ge 4w_-,\\[4pt]
0, & \frac{1}{4}w_- < w_+ < 4w_-,\\[4pt]
2\sqrt{w_+/w_-} - 1, & w_- \ge 4w_+.
\end{cases}
\label{eq:imbalance_func}
\end{equation}
Физический смысл формулы \eqref{eq:3d_reflection} состоит в том, что в режиме сильного дисбаланса ($w_+ \ge 4w_-$) и слабой турбулентности ($\Gamma_+ > \mathcal{R}_{\rm imb}$) скорость передачи амплитуды от выходящей волны к отражённой определяется классическим градиентным отражением Хайнеманна--Олберта $\mathcal{R}_{\rm imb}$ \cite{HeinemannOlbert1980}, тогда как в режиме развитого каскада она не может превышать скорость нелинейной релаксации $\max(\Gamma_+,\Gamma_-)$ и полностью выключается ($\mathcal{C}=0$), когда отношение плотностей энергии достигает порога $w_-/w_+ = 1/4$ (амплитудное отношение $1/2$).

\subsection{Геометрия магнитной трубки и переход к строго консервативной форме энергии}

Рассмотрим течение плазмы вдоль тонкой осесимметричной магнитной силовой трубки с заданной стационарной площадью поперечного сечения $A(r) > 0$, боковая поверхность $\Sigma$ которой образована магнитными силовыми линиями:
\[
\left.\bm{B}\cdot\bm{n}\right|_\Sigma = 0, \qquad \left.\bm{u}\cdot\bm{n}\right|_\Sigma = 0, \qquad \bm{u} \parallel \bm{B}, \qquad \nabla\times\bm{u} = \bm{0}.
\]
Для моделирования как сферически-симметричного расширения, так и сверхрадиальной расходимости корональных дыр на расстояниях $r \sim 1.1\text{--}2.5 R_\odot$ площадь сечения обычно параметризуется через фактор сверхрадиального расширения $f(r)$ Коппа--Хольцера \cite{KoppHolzer1976}:
\begin{equation}
A(r) = A(R_\odot)\left(\frac{r}{R_\odot}\right)^2 f(r), \qquad
f(r) = \frac{f_{\max} + \exp\left(\frac{r_{\rm exp} - r}{\sigma_{\rm exp}}\right)}{1 + \exp\left(\frac{r_{\rm exp} - r}{\sigma_{\rm exp}}\right)}\cdot\frac{1 + \exp\left(\frac{r_{\rm exp} - R_\odot}{\sigma_{\rm exp}}\right)}{f_{\max} + \exp\left(\frac{r_{\rm exp} - R_\odot}{\sigma_{\rm exp}}\right)}.
\label{eq:kopp_holzer}
\end{equation}
Интегрирование условия $\nabla\cdot\bm{B} = 0$ по объёму элемента трубки даёт закон сохранения магнитного потока:
\begin{equation}
A(r)B(r) = \Phi_B = \mathrm{const} \implies V_A(r,\rho) = \frac{\Phi_B}{\sqrt{\mu_0}\,A(r)\sqrt{\rho}}, \qquad L_\perp(r) = C_L\Phi_B^{-1/2}A(r)^{1/2}.
\label{eq:flux_tube}
\end{equation}
Поскольку в стационарном геометрическом русле магнитное поле $B(r) = \Phi_B/A(r)$ не зависит от времени ($\partial B/\partial t = 0$), магнитная энергия $B^2/(2\mu_0)$ и сила магнитного натяжения вдоль линии тока тождественно исключаются из продольных уравнений импульса и энергии.

Введём газодинамическую плотность энергии плазмы $E_g = \frac{1}{2}\rho u^2 + \frac{p}{\gamma-1}$. Прямое усреднение уравнений \eqref{eq:3d_energy}--\eqref{eq:3d_waves} по сечению трубки $A(r)$ приводит к системе с неконсервативным членом работы волнового давления $-u\,\partial P_A/\partial r$ и жёстким источником нагрева $Q_w$:
\begin{align}
\frac{\partial E_g}{\partial t} + \frac{1}{A}\frac{\partial}{\partial r}\left[A u (E_g + p) + A q\right] &= -\rho u \frac{GM_\odot}{r^2} - u\frac{\partial P_A}{\partial r} + Q_w - Q_{\rm rad},
\label{eq:noncons_Eg}\\[4pt]
\frac{\partial w_\pm}{\partial t} + \frac{1}{A}\frac{\partial}{\partial r}\left[A(u \pm V_A)w_\pm\right] + \frac{w_\pm}{2A}\frac{\partial(Au)}{\partial r} &= \mp\mathcal{R}\sqrt{w_+ w_-} - \Gamma_\pm w_\pm.
\label{eq:noncons_w}
\end{align}

Численное интегрирование уравнения \eqref{eq:noncons_Eg} непосредственно для $E_g$ нежелательно по двум причинам: во-первых, аппроксимация неконсервативного произведения $u\,\partial P_A/\partial r$ на крутых градиентах в переходном слое нарушает дискретный баланс механической и волновой энергии; во-вторых, источник $Q_w$ является жёстким при малых временах каскада. Чтобы устранить оба недостатка, заметим, что сумма неконсервативных членов в \eqref{eq:noncons_Eg} и \eqref{eq:noncons_w} образует полную дивергенцию потока работы волнового давления:
\[
u\frac{\partial P_A}{\partial r} + \frac{w_+ + w_-}{2A}\frac{\partial(Au)}{\partial r} = \frac{1}{A}\frac{\partial(A u P_A)}{\partial r},
\]
а сумма источниковых членов диссипации волн и нагрева газа тождественно равна нулю:
\[
-(\Gamma_+ w_+ + \Gamma_- w_-) + Q_w \equiv 0.
\]
Складывая \eqref{eq:noncons_Eg} с обоими уравнениями \eqref{eq:noncons_w}, перейдём к \textbf{полной плотности энергии системы «плазма + альфвеновские волны»} $E$:
\begin{equation}
E = E_g + w_+ + w_- = \frac{\rho u^2}{2} + \frac{p}{\gamma-1} + w_+ + w_-.
\label{eq:total_energy_def}
\end{equation}
Уравнение для $E$ имеет строго дивергентный вид, свободный от $\partial P_A/\partial r$ и $Q_w$:
\begin{equation}
\frac{\partial(A E)}{\partial t} + \frac{\partial}{\partial r}\left[A u(E + p + P_A) + A V_A(w_+ - w_-)\right] = -A\rho u\frac{GM_\odot}{r^2} - \frac{\partial(A q)}{\partial r} - A Q_{\rm rad}.
\label{eq:total_energy_eq_1d}
\end{equation}

\subsection{Модель переходного слоя и регуляризация Линкера--Лионелло--Соколова}

В нижней короне и переходном слое ($T \in [2\times 10^4,\, 2\times 10^6]$~К) перенос энергии вдоль магнитной трубки контролируется балансом классической электронной теплопроводности Спитцера--Харма \cite{SpitzerHarm1953} и оптически тонких радиационных потерь \cite{Rosner1978,Klimchuk2008}:
\begin{equation}
q = -f_c(r)\kappa_0 T^{5/2}\frac{\partial T}{\partial r}, \qquad
Q_{\rm rad}(\rho, T) = n_e n_H \Lambda(T) = c_{\rm rad}\left(\frac{\rho}{m_u}\right)^2\Lambda(T),
\label{eq:tr_physics_raw}
\end{equation}
где $\kappa_0 \approx 9\times 10^{-12}\text{ Вт}\cdot\text{м}^{-1}\cdot\text{К}^{-7/2}$, $c_{\rm rad} \approx 0.61$, а фактор $f_c(r) = \left[1 + (r/r_{\rm col})^2\right]^{-1}$ ($r_{\rm col} \approx 5 R_\odot$) обеспечивает плавное выключение диффузионного закона Фурье при переходе в бесстолкновительный режим на больших гелиоцентрических расстояниях \cite{Hollweg1978}.

В основании переходного слоя ($T \sim 3\times 10^4$~К) малая теплопроводность $\kappa(T) \propto T^{5/2}$ в сочетании с пиком радиационного излучения ионов углерода и кислорода формирует чрезвычайно крутой температурный фронт толщиной $\Delta r_{\rm TR} \sim \frac{\kappa_0 T^{7/2}}{|q_{\rm cor}|} \sim 1\text{--}5$~км ($\sim 10^{-6} R_\odot$). На любой практически реализуемой глобальной сетке ($\Delta r_{\min} \sim 10^{-3} R_\odot \approx 700$~км) такой фронт оказывается недоразрешённым на три порядка величины, что приводит к искусственному подавлению потока тепла из короны в хромосферу и занижению корональной плотности на порядок \cite{Lionello2009}.

Для точного воспроизведения хромосферного испарения на грубой сетке используется аналитическая регуляризация Линкера--Лионелло--Соколова \cite{Lionello2009,Sokolov2013}. Рассмотрим квазистационарный баланс тепла внутри узкого переходного слоя ($p \approx p_{\rm cor} = \mathrm{const}$):
\[
\frac{d}{dr}\left(\kappa(T)\frac{dT}{dr}\right) = Q_{\rm rad}(p_{\rm cor}, T) \implies \frac{1}{2}q_{\rm cor}^2 = \int_{T_0}^{T_{\rm cor}}\kappa(T)Q_{\rm rad}(p_{\rm cor}, T)\,dT.
\]
Поскольку под знаком интеграла стоит исключительно произведение $\kappa(T)\Lambda(T)$, одновременное масштабирование теплопроводности и функции охлаждения при температурах ниже пороговой $T_{\rm mod} \approx 3\times 10^5$~К с коэффициентом $s(T) = \max\left\{1, (T_{\rm mod}/T)^{5/2}\right\}$:
\begin{equation}
\tilde{\kappa}(T) = \kappa_0 T^{5/2} s(T) = \kappa_0\max(T, T_{\rm mod})^{5/2}, \qquad
\tilde{\Lambda}(T) = \frac{\Lambda(T)}{s(T)},
\label{eq:lls_broadening}
\end{equation}
сохраняет тождество $\tilde{\kappa}(T)\tilde{\Lambda}(T) \equiv \kappa(T)\Lambda(T)$. Следовательно, полный интеграл радиационных потерь переходного слоя, давление на вершине слоя $p_{\rm cor}$ и массовый расход солнечного ветра $\dot{m}$ остаются \textbf{строго инвариантными}, тогда как эффективная толщина теплового фронта при $T < T_{\rm mod}$ увеличивается в $(T_{\rm mod}/T_0)^{5/2} \sim 10^2$ раз и гладко разрешается 5--10 ячейками сетки.

\subsection{Векторная запись одномерной системы в консервативных переменных}

Объединяя уравнения сохранения массы, продольного импульса, полной энергии \eqref{eq:total_energy_eq_1d} и волновых плотностей энергии \eqref{eq:noncons_w}, запишем систему в стандартной векторной балансной форме для вектора консервативных переменных $\bm{U}$:
\begin{equation}
\frac{\partial\bm{U}}{\partial t} + \frac{\partial\bm{F}(\bm{U}, r)}{\partial r} + \bm{G}(\bm{U})\frac{\partial(Au)}{\partial r} = \bm{S}_{\rm geom}(\bm{U}, r) + \bm{S}_{\rm turb}(\bm{U}, \partial_r\rho, r) + \bm{S}_{\rm th}(\bm{U}, \partial_r T, \partial_r^2 T, r),
\label{eq:master_system}
\end{equation}
где
\begin{equation}
\bm{U} = A\begin{pmatrix}
\rho \\[3pt]
\rho u \\[3pt]
E \\[3pt]
w_+ \\[3pt]
w_-
\end{pmatrix}, \qquad
\bm{F}(\bm{U}, r) = A\begin{pmatrix}
\rho u \\[3pt]
\rho u^2 + p + P_A \\[3pt]
u(E + p + P_A) + V_A(w_+ - w_-) \\[3pt]
(u + V_A)w_+ \\[3pt]
(u - V_A)w_-
\end{pmatrix}, \qquad
\bm{G}(\bm{U}) = \begin{pmatrix}
0 \\[3pt]
0 \\[3pt]
0 \\[3pt]
\frac{1}{2}w_+ \\[3pt]
\frac{1}{2}w_-
\end{pmatrix},
\label{eq:cons_vectors}
\end{equation}
\begin{equation}
\bm{S}_{\rm geom} = \begin{pmatrix}
0 \\[3pt]
(p + P_A)\dfrac{dA}{dr} - A\rho\dfrac{GM_\odot}{r^2} \\[6pt]
-A\rho u\dfrac{GM_\odot}{r^2} \\[6pt]
0 \\[3pt]
0
\end{pmatrix}, \quad
\bm{S}_{\rm turb} = A\begin{pmatrix}
0 \\[3pt]
0 \\[3pt]
0 \\[3pt]
-\mathcal{R}\sqrt{w_+ w_-} - \Gamma_+ w_+ \\[3pt]
+\mathcal{R}\sqrt{w_+ w_-} - \Gamma_- w_-
\end{pmatrix}, \quad
\bm{S}_{\rm th} = \begin{pmatrix}
0 \\[3pt]
0 \\[3pt]
-\dfrac{\partial(Aq)}{\partial r} - A\tilde{Q}_{\rm rad} \\[6pt]
0 \\[3pt]
0
\end{pmatrix}.
\label{eq:source_vectors}
\end{equation}

% ==============================================================================
\section{Обезразмеривание и параметрический анализ}
% ==============================================================================

\subsection{Анализ выбора характерного масштаба скорости}

На внутренней границе расчётной области $r = R_\odot$ (вершина хромосферы) задаются размерные плотность $\rho_0$, температура $T_0$, магнитное поле $B_0$, поперечный масштаб корреляции $L_{\perp,0}$ и плотность энергии выходящих волн $w_{+,0}^{\rm dim}$. В задаче присутствуют три физических масштаба скорости:
\begin{enumerate}
\item Изотермическая тепловая скорость у основания $c_{T,0} = \sqrt{k_{\rm B}T_0/(\mu m_u)} \approx 20\text{--}30$~км/с (при $T_0 \sim 5\times 10^4$~К);
\item Первая космическая (кеплеровская) скорость на поверхности Солнца: $V_* = v_{\rm K} = \sqrt{GM_\odot/R_\odot} \approx 436.7$~км/с;
\item Альфвеновская скорость у основания $V_{A,0}^{\rm dim} = B_0/\sqrt{\mu_0\rho_0} \approx 200\text{--}1000$~км/с.
\end{enumerate}
Выбор кеплеровской скорости $V_* = \sqrt{GM_\odot/R_\odot}$ в качестве единого масштаба скорости является математически оптимальным по трём причинам: (а) $V_*$ опирается исключительно на фундаментальные константы Солнца $(G, M_\odot, R_\odot)$ и остаётся инвариантной при параметрическом варьировании граничных условий $T_0$, $\rho_0$, $B_0$; (б) гравитационный член в уравнениях импульса и энергии приводится к безразмерному виду $-\rho/r^2$ без дополнительных параметров; (в) в основной области развитого коронального течения ($r > 1.2 R_\odot$) и скорость солнечного ветра ($u \sim 400\text{--}800$~км/с), и корональная скорость звука ($c_s \sim 150\text{--}200$~км/с), и скорость Альфвена имеют порядок $O(V_*)$, что обеспечивает идеальную численную обусловленность матрицы Якоби.

При выборе базисных масштабов
\begin{equation}
\begin{gathered}
R_* = R_\odot, \qquad \rho_* = \rho_0, \qquad V_* = \sqrt{\frac{GM_\odot}{R_\odot}}, \qquad t_* = \frac{R_*}{V_*}, \\
T_* = T_0, \qquad A_* = A(R_\odot), \qquad p_* = \rho_* V_*^2,
\end{gathered}
\label{eq:scales_base}
\end{equation}
безразмерные уравнения (далее записываемые без крышек) принимают вид:
\begin{align}
\frac{\partial(A\rho)}{\partial t} + \frac{\partial(A\rho u)}{\partial r} &= 0,
\label{eq:nd_mass}\\[4pt]
\frac{\partial(A\rho u)}{\partial t} + \frac{\partial}{\partial r}\left[A(\rho u^2 + p + P_A)\right] &= (p + P_A)\frac{dA}{dr} - \frac{A\rho}{r^2},
\label{eq:nd_mom}\\[4pt]
\frac{\partial(A E)}{\partial t} + \frac{\partial}{\partial r}\left[A u(E + p + P_A) + A V_A(w_+ - w_-)\right] &= -\frac{A\rho u}{r^2} - \frac{\partial(A q)}{\partial r} - A\tilde{Q}_{\rm rad},
\label{eq:nd_energy}\\[4pt]
\frac{\partial(A w_\pm)}{\partial t} + \frac{\partial}{\partial r}\left[A(u \pm V_A)w_\pm\right] + \frac{w_\pm}{2}\frac{\partial(Au)}{\partial r} &= A\left(\mp\mathcal{R}\sqrt{w_+ w_-} - \Gamma_\pm w_\pm\right),
\label{eq:nd_waves}
\end{align}
с безразмерными замыканиями:
\begin{equation}
p = \theta_0\rho T, \qquad
V_A = \frac{V_{A,0}}{A\sqrt{\rho}}, \qquad
\Gamma_\pm = \alpha_d\sqrt{\frac{w_\mp}{A\rho}}, \qquad
q = -\chi_0 f_c(r)\max(T, \tau_{\rm mod})^{5/2}\frac{\partial T}{\partial r},
\label{eq:dimless_closures}
\end{equation}
\begin{equation}
\begin{aligned}
\mathcal{R} &= \mathcal{C}(w_+, w_-)\min\left\{V_{A,0}\left|\frac{\partial}{\partial r}\left(\frac{1}{A\sqrt{\rho}}\right)\right|,\, \frac{\alpha_d}{\sqrt{A\rho}}\max\left(\sqrt{w_+}, \sqrt{w_-}\right)\right\}, \\[4pt]
\tilde{Q}_{\rm rad} &= \lambda_{\rm rad}\rho^2 \frac{\hat{\Lambda}(T)}{\max(1, (\tau_{\rm mod}/T)^{5/2})}.
\end{aligned}
\label{eq:dimless_R}
\end{equation}

\subsection{Система безразмерных критериев подобия}

Полный набор управляющих безразмерных параметров $(\theta_0, V_{A,0}, \alpha_d, w_{+,0}, \chi_0, \lambda_{\rm rad}, \tau_{\rm mod})$ и их физическая интерпретация представлены в таблице~\ref{tab:params}.

\begin{table}[ht!]
\centering
\small
\setlength{\tabcolsep}{4pt}
\caption{Безразмерные комплексы подобия модели AWSoM-1D ($R_\odot = 6.96\times 10^8$~м, $V_* = 436.7$~км/с, $t_* = 1593$~с, $\Lambda_* = 10^{-34}\text{ Вт}\cdot\text{м}^3$).}
\vspace{2mm}
\begin{tabular}{@{}llll@{}}
\toprule
\textbf{Параметр} & \textbf{Формула} & \textbf{Физический смысл} & \textbf{Типичные значения} \\
\midrule
$\theta_0$ & $\dfrac{k_{\rm B}T_0}{\mu m_u V_*^2}$ & Отношение хромосферной шкалы высот к $R_\odot$ & $0.0015 \text{ -- } 0.005$ \\[7pt]
$V_{A,0}$ & $\dfrac{B_0}{V_*\sqrt{\mu_0\rho_0}}$ & Безразмерная скорость Альфвена при $r=1$ & $0.5 \text{ -- } 3.0$ \\[7pt]
$\alpha_d$ & $\dfrac{2 R_\odot}{L_{\perp,0}}$ & Безразмерный коэффициент каскадной диссипации & $10 \text{ -- } 100$ \\[7pt]
$w_{+,0}$ & $\dfrac{w_{+,0}^{\rm dim}}{\rho_0 V_*^2}$ & Безразмерная плотность энергии входящих волн & $0.02 \text{ -- } 0.20$ \\[7pt]
$\chi_0$ & $\dfrac{\kappa_0 T_0^{7/2}}{R_\odot\rho_0 V_*^3}$ & Обратное число Пекле теплопроводности при $T=T_0$ & $10^{-7} \text{ -- } 10^{-5}$ \\[7pt]
$\lambda_{\rm rad}$ & $\dfrac{c_{\rm rad} R_\odot \rho_0 \Lambda_*}{m_u^2 V_*^3}$ & Безразмерный параметр радиационных потерь & $0.05 \text{ -- } 2.0$ \\[7pt]
$\tau_{\rm mod}$ & $\dfrac{T_{\rm mod}}{T_0}$ & Порог уширения переходного слоя & $5.0 \text{ -- } 15.0$ \\
\bottomrule
\end{tabular}
\label{tab:params}
\end{table}

% ==============================================================================
\section{Математический анализ системы уравнений}
% ==============================================================================

\subsection{Скрытая квазилинейность коэффициента отражения и три ветви системы}

Принципиальная математическая особенность системы \eqref{eq:nd_mass}--\eqref{eq:nd_waves} с замыканием \eqref{eq:dimless_R} заключается в структуре градиентной ветви коэффициента отражения $\mathcal{R}$. Раскрывая производную от скорости Альфвена $V_A(r, \rho) = V_{A,0} A(r)^{-1}\rho^{-1/2}$ по радиальной координате:
\begin{equation}
\frac{\partial V_A}{\partial r} = -V_A\left(\frac{A'}{A} + \frac{1}{2\rho}\frac{\partial\rho}{\partial r}\right),
\label{eq:dVA_dr}
\end{equation}
обнаруживаем, что величина $|\partial V_A/\partial r|$ явно зависит от первой пространственной производной плотности $\partial\rho/\partial r$. Это влечёт за собой два фундаментальных математических следствия:
\begin{enumerate}
\item \textbf{В нестационарной задаче:} когда в операторе $\min\{\dots\}$ формулы \eqref{eq:dimless_R} активна первая (градиентная) ветвь, слагаемое $\mp\mathcal{R}\sqrt{w_+ w_-}$ в правой части волновых уравнений \eqref{eq:nd_waves} перестаёт быть алгебраическим источником нулевого порядка и становится \textbf{неконсервативным дифференциальным членом первого порядка} $\propto \frac{\partial\rho}{\partial r}$, который должен быть перенесён в левую часть и включён в главную матрицу Якоби системы.
\item \textbf{В стационарной редукции ($\partial/\partial t = 0$):} в силу интеграла сохранения массы $\rho = \dot{m}/(Au)$ производная плотности выражается через производную скорости: $-\frac{\rho'}{\rho} = \frac{u'}{u} + \frac{A'}{A}$, откуда $|\frac{dV_A}{dr}| = \frac{V_A}{2}\left|\frac{u'}{u} - \frac{A'}{A}\right|$. Следовательно, стационарная система представляет собой \textbf{неявное кусочно-гладкое ОДУ (систему типа Филиппова), не разрешённое относительно старшей производной $u'$}, которое в каждой точке распадается на три аналитические ветви (диссипативную ветвь и две градиентные ветви противоположного знака $\frac{u'}{u} - \frac{A'}{A} \gtrless 0$).
\end{enumerate}

\subsection{Квазилинейная форма и характеристический анализ (гиперболичность)}

Исследуем конвективно-волновую часть нестационарной системы \eqref{eq:nd_mass}--\eqref{eq:nd_waves} (при выделении главного гиперболического оператора первого порядка) в векторе примитивных переменных $\bm{W} = (\rho, u, p, w_+, w_-)^T$. Введём кусочно-постоянный индикатор активной ветви отражения $\sigma$:
\begin{equation}
\sigma =
\begin{cases}
0, & \text{если } \left|\dfrac{\partial V_A}{\partial r}\right| \ge \max(\Gamma_+,\Gamma_-) \quad \text{(диссипативная ветвь)},\\[8pt]
\mathrm{sgn}\left(\dfrac{\partial V_A}{\partial r}\right) = \pm 1, & \text{если } \left|\dfrac{\partial V_A}{\partial r}\right| < \max(\Gamma_+,\Gamma_-) \quad \text{(градиентные ветви)}.
\end{cases}
\label{eq:sigma_branch}
\end{equation}
Используя тождество $\frac{\partial V_A}{\partial r} = -\frac{V_A}{2\rho}\frac{\partial\rho}{\partial r} - V_A\frac{A'}{A}$, приведём систему к квазилинейному виду
\begin{equation}
\frac{\partial\bm{W}}{\partial t} + \mathsf{A}_{\rm prim}(\bm{W})\frac{\partial\bm{W}}{\partial r} = \bm{S}_{\rm alg}(\bm{W}, r) + \bm{S}_{\rm diff}(\bm{W}, \partial_r^2 T),
\label{eq:quasilinear_form}
\end{equation}
где матрица $\mathsf{A}_{\rm prim}(\bm{W})$ размера $5 \times 5$ равна:
\begin{equation}
\mathsf{A}_{\rm prim}(\bm{W}) =
\begin{pmatrix}
u & \rho & 0 & 0 & 0 \\[4pt]
0 & u & \dfrac{1}{\rho} & \dfrac{1}{2\rho} & \dfrac{1}{2\rho} \\[6pt]
0 & \gamma p & u & 0 & 0 \\[4pt]
-\alpha_+ & \dfrac{3}{2}w_+ & 0 & u + V_A & 0 \\[6pt]
\alpha_- & \dfrac{3}{2}w_- & 0 & 0 & u - V_A
\end{pmatrix}, \qquad
\alpha_\pm = \frac{V_A}{2\rho}\left(w_\pm - \sigma\mathcal{C}\sqrt{w_+ w_-}\right).
\label{eq:A_prim}
\end{equation}

Найдём характеристические скорости $\lambda$ из уравнения $\det(\mathsf{A}_{\rm prim} - \lambda\mathbf{I}) = 0$. Вводя собственную скорость возмущений относительно движущейся плазмы $\tilde{\lambda} = \lambda - u$ и раскрывая определитель по первому и третьему столбцам, получаем:
\begin{equation}
\tilde{\lambda}\cdot P_4(\tilde{\lambda}) = 0,
\label{eq:det_factor}
\end{equation}
где $\tilde{\lambda}_1 = 0$ ($\lambda_1 = u$) отвечает контактной (энтропийной) характеристике, а полином четвёртой степени $P_4(\tilde{\lambda})$ описывает зацепление акустических и альфвеновских волн огибающей через пондеромоторное давление и отражение:
\begin{equation}
P_4(\tilde{\lambda}) = \left(\tilde{\lambda}^2 - c_s^2\right)\left(\tilde{\lambda}^2 - V_A^2\right) - \frac{3(w_+ + w_-)}{4\rho}\tilde{\lambda}^2 - \tilde{\lambda}V_A\left[\frac{w_+ - w_-}{2\rho} + \sigma\mathcal{C}\frac{\sqrt{w_+ w_-}}{2\rho}\right] + \frac{w_+ + w_-}{4\rho}V_A^2 = 0,
\label{eq:P4_poly}
\end{equation}
где $c_s^2 = \gamma p/\rho$. Анализ уравнения \eqref{eq:P4_poly} позволяет сформулировать ключевые структурные теоремы о системе:
\begin{enumerate}
\item \textbf{На диссипативной ветви ($\sigma = 0$):} уравнение \eqref{eq:P4_poly} совпадает с классическим характеристическим уравнением МГД с волновым давлением ВКБ \cite{Jacques1977,vanderHolst2014}. Вводя безразмерные малые параметры относительной плотности волновой энергии $\epsilon_\pm = \frac{w_\pm}{\rho V_A^2} \ll 1$, методом возмущений при $c_s \neq V_A$ находим четыре вещественных различных корня:
\begin{align}
\tilde{\lambda}_{2,3} &= \pm c_{\rm eff} + O(\epsilon^2), \qquad c_{\rm eff}^2 = c_s^2 + \frac{w_+ + w_-}{4\rho}\left(\frac{V_A^2 - 3c_s^2}{V_A^2 - c_s^2}\right),
\label{eq:ceff_roots}\\[4pt]
\tilde{\lambda}_{4,5} &= \pm V_A\left(1 + \frac{w_\pm}{2\rho(V_A^2 - c_s^2)}\right) + O(\epsilon^2).
\label{eq:va_roots}
\end{align}
Заметим замечательный физический эффект в формуле \eqref{eq:ceff_roots}: в доальфвеновской короне ($V_A \gg c_s$) эффективный квадрат скорости звука равен $c_{\rm eff}^2 \approx c_s^2 + \frac{w_+ + w_-}{4\rho}$ (эффективный показатель адиабаты волнового газа равен $\gamma_w = 1/2$), тогда как в далёком суперальфвеновском ветре ($c_s^2 \ll V_A^2$ нарушается или при формальном пределе $V_A \to 0$) $c_{\rm eff}^2 \approx c_s^2 + \frac{3(w_+ + w_-)}{4\rho}$ ($\gamma_w = 3/2$, результат Жака \cite{Jacques1977}).
\item \textbf{На градиентной ветви ($\sigma = \pm 1$):} слагаемое с отражением $\sigma\mathcal{C}\sqrt{w_+ w_-}$ сокращается в свободном члене полинома \eqref{eq:P4_poly}, но входит в коэффициент при первой степени $\tilde{\lambda} V_A$, создавая дополнительный асимметричный сдвиг фазовых скоростей порядка $\Delta\tilde{\lambda} \sim \frac{\sigma\mathcal{C}\sqrt{w_+ w_-}}{4\rho V_A}$. Поскольку в солнечной короне и ветре $\frac{\sqrt{w_+ w_-}}{\rho V_A^2} \sim 10^{-3}\text{--}10^{-2} \ll 1$, дискриминант полинома $P_4(\tilde{\lambda})$ остаётся строго положительным. Таким образом, \textbf{конвективно-волновая подсистема является строго гиперболической на всех трёх ветвях замыкания отражения}.
\item \textbf{Разрывность матрицы Якоби на поверхностях переключения:} при пересечении поверхностей смены ветви $|\partial_r V_A| = \max(\Gamma_+,\Gamma_-)$ и $\partial_r V_A = 0$ индикатор $\sigma$ меняется скачком, вследствие чего матрица Якоби $\mathsf{A}_{\rm prim}(\bm{W})$ терпит разрыв первого рода. Именно это обстоятельство делает применение матричных схем расщепления типа Роу (Roe solver), требующих гладкой аналитической диагонализации $\mathsf{A} = \mathsf{R}\bm{\Lambda}\mathsf{R}^{-1}$, неустойчивым и диктует выбор центрально-противопоточных схем типа HLL, использующих лишь внешние границы спектра $[S_L, S_R]$.
\end{enumerate}

\subsection{Первые интегралы и топология критических точек в стационарном пределе}

При выходе нестационарного решения на стационарный режим ($\partial/\partial t \to 0$) уравнения массы \eqref{eq:nd_mass} и полной энергии \eqref{eq:nd_energy} дают два точных первых интеграла вдоль силовой трубки --- постоянный безразмерный массовый расход $\dot{m}$ и обобщённый интеграл Бернулли $\varepsilon^{\rm tot}$:
\begin{equation}
\dot{m} = A(r)\rho(r)u(r) = \mathrm{const},
\label{eq:int_mass}
\end{equation}
\begin{equation}
\varepsilon^{\rm tot} = \frac{u^2}{2} + \frac{\gamma}{\gamma-1}\frac{p}{\rho} - \frac{1}{r} + \frac{3(w_+ + w_-)}{2\rho} + \frac{V_A(w_+ - w_-)}{\rho u} + \frac{A q}{\dot{m}} + \frac{1}{\dot{m}}\int_1^r A(r')\tilde{Q}_{\rm rad}\,dr' = \mathrm{const}.
\label{eq:int_bernoulli}
\end{equation}

В отсутствие теплопроводности ($q=0$) стационарная система сводится к системе трёх ОДУ для $(u, w_+, w_-)$ вида $\mathsf{M}(r, \bm{W})\bm{W}' = \bm{f}(r, \bm{W})$ с определителем $\det\mathsf{M} = (u + V_A)(u - V_A)D(r, \bm{W})$, который обращается в нуль в двух физически различных критических точках, где одна из характеристик $\lambda_k$ меняет знак (переходит через нуль в лабораторной системе координат):
\begin{enumerate}
\item \textbf{Модифицированная звуковая точка ($r = r_s$, где $\lambda_3 = u - c_{\rm eff} = 0$, $D = 0$):} точка перехода течения от дозвукового к сверхзвуковому режиму ($r_s \approx 1.8\text{--}3.5 R_\odot$). Условие регулярности производной $u'$ ($N = 0$ при $D = 0$) однозначно фиксирует безразмерный массовый расход ветра $\dot{m} = u(1)$.
\item \textbf{Альфвеновская критическая точка ($r = r_A$, где $\lambda_5 = u - V_A = 0$):} точка перехода от субальфвеновского к суперальфвеновскому течению ($r_A \approx 8\text{--}18 R_\odot$). В этой точке коэффициент при производной встречной волны $(u - V_A)\frac{dw_-}{dr}$ в уравнении для $w_-$ обращается в нуль! Следовательно, регулярное прохождение решения через точку Альфвена требует зануления правой части третьего уравнения в $r = r_A$, что однозначно фиксирует второй свободный параметр на нижней границе --- плотность энергии приходящих к Солнцу отражённых волн $w_-(1)$.
\end{enumerate}

% ==============================================================================
\section{Вычислительные трудности и архитектура численной схемы}
% ==============================================================================

Прямое решение краевой стационарной задачи методом пристрелки между тремя точками ($r = 1$, звуковой точкой $r_s$ и альфвеновской точкой $r_A$) на фоне переключения трёх негладких ветвей $\mathcal{R}$ и параболического пограничного слоя Спитцера--Харма представляет собой исключительно плохо обусловленную двухпараметрическую задачу. В противоположность этому, \textbf{метод установления по времени (time relaxation)} гиперболо-параболической системы \eqref{eq:nd_mass}--\eqref{eq:nd_waves} позволяет потоку самостоятельно и автоматически сформировать регулярный переход через обе критические точки.

Однако для успешной реализации метода установления численная схема должна преодолеть пять фундаментальных вычислительных трудностей, сводка которых и выбранные методы решения приведены ниже.

\subsection{Систематизация численных трудностей}

\begin{enumerate}
\item \textbf{Экстремальная характеристическая жёсткость (CFL-барьер):}
В нижней короне ($r \approx 1.05\text{--}1.5 R_\odot$) плотность падает на три порядка, вследствие чего скорость Альфвена $V_A \propto \rho^{-1/2}$ достигает $V_A \sim 1500\text{--}3000$~км/с ($\hat{V}_A \sim 4\text{--}7$), тогда как скорость втекания в переходном слое составляет лишь $u \sim 1\text{--}5$~км/с ($\hat{u} \sim 10^{-2}$). В сочетании с мелким шагом сетки у границы ($\Delta r_{\min} \sim 10^{-3}$) явный шаг Куранта $\Delta t_{\rm CFL} \sim \Delta r_{\min}/V_A \sim 2\times 10^{-4}$ требует порядка $10^5$ шагов по времени для полной релаксации медленных акустических и энтропийных мод во всей расчётной области $r \in [1, 30]$.
\emph{Решение:} Геометрически растянутая сетка $\Delta r_{i+1} = q\Delta r_i$ ($q \approx 1.008$), увеличивающая шаг ячейки пропорционально росту характерного масштаба градиентов, и параллелизация всех пространственных циклов с помощью OpenMP.

\item \textbf{Параболическая жёсткость электронной теплопроводности:}
Коэффициент теплопроводности Спитцера $\tilde{\kappa}(T) \propto T^{5/2}$ возрастает от хромосферы к короне в $35^{5/2} \approx 7000$ раз. Явное условие устойчивости параболического оператора $\Delta t_{\rm diff} \le \frac{\Delta r^2}{2\chi_{\rm eff}}$ в переходном слое оказывается в $10^2\text{--}10^3$ раз жёстче гиперболического шага Куранта $\Delta t_{\rm CFL}$.
\emph{Решение:} Выделение оператора теплопроводности и радиационных потерь $\bm{S}_{\rm th}$ в отдельный глобальный неявный подшаг расщепления, решаемый на каждом шаге по времени методом трёхдиагональной прогонки Томаса за $O(N)$ операций при безусловной устойчивости.

\item \textbf{Источниковая жёсткость каскада и корневая сингулярность $\sqrt{w_+ w_-}$:}
В нижней короне безразмерный параметр диссипации велик ($\alpha_d \sim 30\text{--}100$), поэтому локальное время затухания отражённой волны $\tau_- = \Gamma_-^{-1} = \frac{\sqrt{A\rho}}{\alpha_d\sqrt{w_+}}$ оказывается меньше $\Delta t_{\rm CFL}$. Более того, правые части уравнений \eqref{eq:nd_waves} содержат нелипшицев корень $\sqrt{w_+ w_-}$, якобиан которого по $w_-$ расходится как $w_-^{-1/2} \to \infty$ при $w_- \to 0$. В любой стандартной явной схеме Рунге--Кутты это неизбежно приводит к пробиванию нуля ($w_-^{n+1} < 0$) и разрушению расчёта из-за комплексного значения $\sqrt{w_-}$.
\emph{Решение:} Аналитический переход на локальном источниковом шаге от плотностей энергии $w_\pm$ к \textbf{амплитудным переменным} $a_\pm = \sqrt{w_\pm}$, который полностью устраняет корневую сингулярность и сводит систему ОДУ к гладкой билинейной системе $2 \times 2$, решаемой полунеявно с гарантией строгой положительности $w_\pm^{n+1} \ge 0$.

\item \textbf{Неконсервативные члены и согласованность баланса на неоднородной сетке:}
Присутствие геометрического члена $(p + P_A)dA/dr$ в уравнении импульса и неконсервативного сжатия $\frac{w_\pm}{2}\frac{\partial(Au)}{\partial r}$ в уравнениях волн может порождать ошибки аппроксимации на сильно растянутой сетке.
\emph{Решение:} Использование полной консервативной энергии $E = E_g + w_+ + w_-$ (полностью исключающей неконсервативный член $-u\,\partial_r P_A$) и вычисление дивергенции объёмного потока $\partial_r(Au)$ непосредственно через численные потоки массы на гранях ячеек.

\item \textbf{Негладкость операторов $\min/\max$ в коэффициенте отражения $\mathcal{R}$:}
В точках переключения градиентной и диссипативной ветвей производные правой части разрывны, что при подходе к стационару может вызывать высокочастотное дрожание (chattering) невязки на уровне $10^{-5}$.
\emph{Решение:} Возможность переключения на $C^\infty$-гладкую регуляризацию оператора $\min(X, Y)$ через гармоническое среднее $\frac{XY}{\sqrt{X^2 + Y^2}}$ \cite{Cranmer2007}.
\end{enumerate}

\subsection{Пошаговый алгоритм численной схемы (MUSCL-HLL + RK2 + IMEX + Томас)}

Расчётная область $r \in [r_{\rm in}, r_{\rm out}]$ ($r_{\rm in} = 1$, $r_{\rm out} = 30$) покрывается сеткой из $N$ ячеек $[r_{i-1/2}, r_{i+1/2}]$ с центрами $r_i = \frac{1}{2}(r_{i-1/2} + r_{i+1/2})$ и шагами $\Delta r_i = \Delta r_1 q^{i-1}$, дополненной двумя фиктивными ячейками (ghost cells) с каждого края. На каждом временном шаге $\Delta t = C_{\rm CFL}\min_i \frac{\Delta r_i}{|u_i| + c_{{\rm fast},i}}$ выполняются три стадии.

\paragraph{Стадия 1. Пространственный гиперболический оператор (MUSCL-HLL / Rusanov).}
По известным средним значениям примитивных переменных $\bm{W}_i$ выполняется кусочно-линейная реконструкция MUSCL на внутренние стороны граней $r_{i+1/2}$ с учётом неравномерности шагов сетки:
\begin{equation}
\begin{aligned}
\bm{W}_{i+1/2}^L &= \bm{W}_i + \frac{\Delta r_i}{2}\,\mathcal{L}\!\left(\frac{\bm{W}_i - \bm{W}_{i-1}}{r_i - r_{i-1}},\, \frac{\bm{W}_{i+1} - \bm{W}_i}{r_{i+1} - r_i}\right), \\[4pt]
\bm{W}_{i+1/2}^R &= \bm{W}_{i+1} - \frac{\Delta r_{i+1}}{2}\,\mathcal{L}\!\left(\frac{\bm{W}_{i+1} - \bm{W}_i}{r_{i+1} - r_i},\, \frac{\bm{W}_{i+2} - \bm{W}_{i+1}}{r_{i+2} - r_{i+1}}\right),
\end{aligned}
\label{eq:muscl_nonuniform}
\end{equation}
где $\mathcal{L}(a, b) = \frac{2ab}{a+b}$ (при $ab > 0$) --- симметричный монотонный лимитер ван Лира. По реконструированным состояниям $\bm{W}_{i+1/2}^{L,R}$ и точному значению площади на грани $A_{i+1/2} = A(r_{i+1/2})$ вычисляется численный поток HLL \cite{HLL1983}:
\begin{equation}
\hat{\bm{F}}_{i+1/2} =
\begin{cases}
\bm{F}_L, & S_L \ge 0,\\[4pt]
\dfrac{S_R\bm{F}_L - S_L\bm{F}_R + S_L S_R(\bm{U}_R - \bm{U}_L)}{S_R - S_L}, & S_L < 0 < S_R,\\[8pt]
\bm{F}_R, & S_R \le 0,
\end{cases}
\label{eq:hll_formula}
\end{equation}
где крайние скорости волн и быстрая магнитозвуковая мажоранта равны:
\[
S_L = \min(u_L - c_{{\rm fast},L},\, u_R - c_{{\rm fast},R}), \quad
S_R = \max(u_L + c_{{\rm fast},L},\, u_R + c_{{\rm fast},R}), \quad
c_{\rm fast} = \max\left(V_A, \sqrt{c_s^2 + \frac{3(w_+ + w_-)}{4\rho}}\right).
\]
Аналогичной HLL-формулой вычисляется объёмный расход на грани $(Au)_{i+1/2}$ для аппроксимации неконсервативного члена $-\frac{w_{\pm,i}}{2}\frac{(Au)_{i+1/2} - (Au)_{i-1/2}}{\Delta r_i}$. Интегрирование гиперболической части выполняется двухшаговой TVD-схемой Рунге--Кутты второго порядка (SSP-RK2).

\paragraph{Стадия 2. Локальный полунеявный шаг для волновых источников в амплитудах.}
После каждого подшага Рунге--Кутты в каждой ячейке $i$ при фиксированных плотности $\rho_i$, скорости $u_i$ и полной энергии $E_i^*$ интегрируется локальная система ОДУ:
\begin{equation}
\frac{dw_+}{dt} = -\mathcal{R}_i\sqrt{w_+ w_-} - k_{d,i} w_+\sqrt{w_-}, \qquad
\frac{dw_-}{dt} = +\mathcal{R}_i\sqrt{w_+ w_-} - k_{d,i} w_-\sqrt{w_+}, \qquad k_{d,i} = \frac{\alpha_d}{\sqrt{A_i\rho_i}}.
\label{eq:local_wave_ode}
\end{equation}
Подстановка $w_+ = a_+^2$, $w_- = a_-^2$ ($dw_\pm = 2a_\pm da_\pm$) позволяет сократить сингулярные множители $a_\pm$ и записать линейную относительно $(a_+^{**}, a_-^{**})$ полунеявную схему:
\begin{equation}
\begin{pmatrix}
1 & \dfrac{\Delta t}{2}\left(\mathcal{R}_i + k_{d,i} a_+^*\right) \\[8pt]
-\dfrac{\Delta t}{2}\mathcal{R}_i & 1 + \dfrac{\Delta t}{2}k_{d,i} a_+^*
\end{pmatrix}
\begin{pmatrix}
a_+^{**} \\[6pt]
a_-^{**}
\end{pmatrix} =
\begin{pmatrix}
a_+^* \\[6pt]
a_-^*
\end{pmatrix}.
\label{eq:imex_2x2}
\end{equation}
Определитель системы \eqref{eq:imex_2x2} равен $\Delta = 1 + \frac{\Delta t}{2}k_{d,i}a_+^* + \frac{\Delta t^2}{4}\mathcal{R}_i(\mathcal{R}_i + k_{d,i}a_+^*) > 1$ при типичном положительном коэффициенте отражения $\mathcal{R}_i \ge 0$, что гарантирует безусловную устойчивость и неотрицательность $w_\pm^{**} = (a_\pm^{**})^2 \ge 0$. Поскольку полная энергия $U_{3,i} = A_i E_i^*$ на этом шаге не меняется, новое тепловое давление восстанавливается из закона сохранения энергии:
\[
p_i^{**} = (\gamma - 1)\left(E_i^* - \frac{1}{2}\rho_i u_i^2 - w_{+,i}^{**} - w_{-,i}^{**}\right).
\]

\paragraph{Стадия 3. Неявная трёхдиагональная прогонка Томаса для переходного слоя.}
На заключительном этапе шага $\Delta t$ решается неявное разностное уравнение теплопроводности и радиационных потерь для безразмерной температуры $T_i^{n+1}$ ($i = 1, \dots, N$):
\begin{equation}
\begin{aligned}
\frac{\theta_0}{\gamma-1}A_i\rho_i\frac{T_i^{n+1} - T_i^{**}}{\Delta t} &= \frac{1}{\Delta r_i}\left(K_{i+1/2}\frac{T_{i+1}^{n+1} - T_i^{n+1}}{r_{i+1} - r_i} - K_{i-1/2}\frac{T_i^{n+1} - T_{i-1}^{n+1}}{r_i - r_{i-1}}\right) \\[4pt]
&\quad - A_i\left[\frac{\tilde{Q}_{\rm rad}(\rho_i, T_i^{**})}{T_i^{**}}\right]T_i^{n+1},
\end{aligned}
\label{eq:thomas_eq}
\end{equation}
где $K_{i+1/2} = A_{i+1/2} f_c(r_{i+1/2})\chi_0\max\!\left(\frac{T_i^{**} + T_{i+1}^{**}}{2}, \tau_{\rm mod}\right)^{5/2}$. Матрица системы \eqref{eq:thomas_eq} является трёхдиагональной матрицей со строгим диагональным преобладанием ($b_i > a_i + c_i$, $a_i, c_i > 0$), поэтому алгоритм прогонки Томаса находит единственное строго положительное распределение температуры $T_i^{n+1} > 0$ за $O(N)$ арифметических операций.

\subsection{Характеристические граничные условия на входе и выходе}

\begin{itemize}
\item \textbf{Внутренняя граница ($r = 1$, вершина хромосферы):} так как течение у основания дозвуковое и субальфвеновское ($0 < u < c_{\rm eff} < V_A$), три характеристики входят в область ($\lambda_1 = u > 0$, $\lambda_2 = u + c_{\rm eff} > 0$, $\lambda_4 = u + V_A > 0$), а две выходят из области влево ($\lambda_3 = u - c_{\rm eff} < 0$, $\lambda_5 = u - V_A < 0$). Поэтому на границе $r = 1$ жёстко фиксируются ровно три параметра входящих волн:
\[
\rho(1, t) = 1, \qquad p(1, t) = \theta_0 \quad (T(1,t) = 1), \qquad w_+(1, t) = w_{+,0},
\]
тогда как скорость втекания $u(1, t)$ и энергия приходящих отражённых волн $w_-(1, t)$ сносятся из первой расчётной ячейки по выходящим характеристикам ($u_0 = u_1$, $w_{-,0} = w_{-,1}$).
\item \textbf{Внешняя граница ($r = r_{\rm out} = 30$, сверхзвуковой и суперальфвеновский ветер):} поскольку $u > V_A > c_{\rm eff}$, все пять характеристических скоростей строго положительны ($\lambda_{1..5} > 0$), и на правой границе ставятся неотражающие условия свободной экстраполяции $\bm{W}_{N+1} = \bm{W}_N$.
\end{itemize}

% ==============================================================================
\begin{thebibliography}{99}

\bibitem{Parker1958}
Parker, E.~N. Dynamics of the Interplanetary Gas and Magnetic Fields // \emph{The Astrophysical Journal}. --- 1958. --- Vol.~128. --- P.~664--676.

\bibitem{BelcherDavis1971}
Belcher, J.~W., Davis, L., Jr. Large-amplitude Alfv\'en waves in the interplanetary medium, 2 // \emph{Journal of Geophysical Research}. --- 1971. --- Vol.~76, No.~16. --- P.~3534--3563.

\bibitem{Alazraki1971}
Alazraki, G., Couturier, P. Solar wind acceleration caused by the gradient of Alfv\'en wave pressure // \emph{Astronomy and Astrophysics}. --- 1971. --- Vol.~13. --- P.~380--389.

\bibitem{Belcher1971}
Belcher, J.~W. Alfv\'enic wave pressures and the solar wind // \emph{The Astrophysical Journal}. --- 1971. --- Vol.~168. --- P.~509--524.

\bibitem{Hollweg1973}
Hollweg, J.~V. Alfv\'en waves in a two-fluid model of the solar wind // \emph{The Astrophysical Journal}. --- 1973. --- Vol.~181. --- P.~547--566.

\bibitem{Hollweg1978}
Hollweg, J.~V. Some physical processes in the solar wind // \emph{Reviews of Geophysics and Space Physics}. --- 1978. --- Vol.~16, No.~4. --- P.~689--720.

\bibitem{Jacques1977}
Jacques, S.~A. Momentum and energy transport by waves in the solar atmosphere and solar wind // \emph{The Astrophysical Journal}. --- 1977. --- Vol.~215. --- P.~942--951.

\bibitem{Kraichnan1965}
Kraichnan, R.~H. Inertial-range spectrum of hydromagnetic turbulence // \emph{Physics of Fluids}. --- 1965. --- Vol.~8, No.~7. --- P.~1385--1387.

\bibitem{Dobrowolny1980}
Dobrowolny, M., Mangeney, A., Veltri, P. Fully developed anisotropic hydromagnetic turbulence in interplanetary space // \emph{Physical Review Letters}. --- 1980. --- Vol.~45, No.~2. --- P.~144--147.

\bibitem{HeinemannOlbert1980}
Heinemann, M., Olbert, S. Non-WKB Alfv\'en waves in the solar wind // \emph{Journal of Geophysical Research}. --- 1980. --- Vol.~85, No.~A3. --- P.~1311--1327.

\bibitem{Velli1993}
Velli, M. On the propagation of ideal, linear Alfv\'en waves in radially stratified stellar atmospheres and winds // \emph{Astronomy and Astrophysics}. --- 1993. --- Vol.~270. --- P.~304--314.

\bibitem{Matthaeus1999}
Matthaeus, W.~H., Zank, G.~P., Oughton, S., Mullan, D.~J., Dmitruk, P. Coronal heating by magnetohydrodynamic turbulence driven by reflected low-frequency waves // \emph{The Astrophysical Journal Letters}. --- 1999. --- Vol.~523, No.~1. --- P.~L93--L96.

\bibitem{GoldreichSridhar1995}
Goldreich, P., Sridhar, S. Toward a theory of interstellar turbulence. II. Strong alfvenic turbulence // \emph{The Astrophysical Journal}. --- 1995. --- Vol.~438. --- P.~763--775.

\bibitem{Dmitruk2002}
Dmitruk, P., Matthaeus, W.~H., Milano, L.~J., Oughton, S., Zank, G.~P., Mullan, D.~J. Coronal heating distribution due to low-frequency, wave-driven turbulence // \emph{The Astrophysical Journal}. --- 2002. --- Vol.~575, No.~1. --- P.~571--577.

\bibitem{SuzukiInutsuka2005}
Suzuki, T.~K., Inutsuka, S.-i. Making the corona and the fast solar wind: a self-consistent simulation for the low-frequency Alfv\'en waves from the photosphere to 0.3 AU // \emph{The Astrophysical Journal Letters}. --- 2005. --- Vol.~632, No.~1. --- P.~L49--L52.

\bibitem{Cranmer2007}
Cranmer, S.~R., van Ballegooijen, A.~A., Edgar, R.~J. Self-Consistent Coronal Heating and Solar Wind Acceleration from Anisotropic Magnetohydrodynamic Turbulence // \emph{The Astrophysical Journal Supplement Series}. --- 2007. --- Vol.~171, No.~2. --- P.~520--551.

\bibitem{VerdiniVelli2007}
Verdini, A., Velli, M. Alfv\'en waves and turbulence in the solar atmosphere and solar wind // \emph{The Astrophysical Journal}. --- 2007. --- Vol.~662, No.~1. --- P.~669--676.

\bibitem{Verdini2010}
Verdini, A., Velli, M., Matthaeus, W.~H., Oughton, S., Dmitruk, P. A Turbulence-Driven Model for Heating and Acceleration of the Fast Solar Wind in Coronal Holes // \emph{The Astrophysical Journal Letters}. --- 2010. --- Vol.~708, No.~2. --- P.~L116--L120.

\bibitem{Chandran2009}
Chandran, B.~D.~G., Hollweg, J.~V. Alfv\'en wave reflection and turbulent heating in the solar wind from 1 solar radius to 1 AU: An analytical treatment // \emph{The Astrophysical Journal}. --- 2009. --- Vol.~707, No.~2. --- P.~1659--1667.

\bibitem{Chandran2011}
Chandran, B.~D.~G., Dennis, T.~J., Quataert, E., Bale, S.~D. Incorporating Low-Frequency Alfv\'en-Wave Turbulence into a One-Dimensional Solar-Wind Model // \emph{The Astrophysical Journal}. --- 2011. --- Vol.~743, No.~2. --- Art.~197.

\bibitem{Hollweg1986}
Hollweg, J.~V. Transition region, corona, and solar wind in coronal holes // \emph{Journal of Geophysical Research}. --- 1986. --- Vol.~91, No.~A4. --- P.~4111--4125.

\bibitem{Sokolov2013}
Sokolov, I.~V., van der Holst, B., Oran, R., Downs, C., Roussev, I.~I., Jin, M., Manchester, W.~B., IV, Evans, R.~M., Gombosi, T.~I. Magnetohydrodynamic Waves and Coronal Heating: Unifying Empirical and MHD Turbulence Models // \emph{The Astrophysical Journal}. --- 2013. --- Vol.~764, No.~1. --- Art.~23.

\bibitem{Oran2013}
Oran, R., van der Holst, B., Landi, E., Jin, M., Sokolov, I.~V., Gombosi, T.~I. A Global Wave-Driven Magnetohydrodynamic Solar Model with a Unified Treatment of Open and Closed Magnetic Field Topologies // \emph{The Astrophysical Journal}. --- 2013. --- Vol.~778, No.~2. --- Art.~176.

\bibitem{vanderHolst2014}
van der Holst, B., Sokolov, I.~V., Meng, X., Jin, M., Manchester, W.~B., IV, T\'oth, G., Gombosi, T.~I. Alfv\'en Wave Solar Model (AWSoM): Coronal Heating // \emph{The Astrophysical Journal}. --- 2014. --- Vol.~782, No.~2. --- Art.~81.

\bibitem{Sokolov2016}
Sokolov, I.~V., van der Holst, B., Manchester, W.~B., Su Ozturk, D., Szente, J., Taktakishvili, A., T\'oth, G., Jin, M., Gombosi, T.~I. Threaded-Field-Line Model for the Low Solar Corona Powered by the Alfv\'en Wave Turbulence // \emph{arXiv preprint arXiv:1609.04379}. --- 2016.

\bibitem{Sokolov2021}
Sokolov, I.~V., van der Holst, B., Manchester, W.~B., Su Ozturk, D., Szente, J., Taktakishvili, A., T\'oth, G., Jin, M., Gombosi, T.~I. Threaded-Field-Line Model for the Low Solar Corona Powered by the Alfv\'en Wave Turbulence // \emph{The Astrophysical Journal}. --- 2021. --- Vol.~908, No.~2. --- Art.~172.

\bibitem{KoppHolzer1976}
Kopp, R.~A., Holzer, T.~E. Dynamics of coronal hole regions. I. Steady polytropic flows with multiple critical points // \emph{Solar Physics}. --- 1976. --- Vol.~49, No.~1. --- P.~43--56.

\bibitem{SpitzerHarm1953}
Spitzer, L., H\"arm, R. Transport phenomena in a completely ionized gas // \emph{Physical Review}. --- 1953. --- Vol.~89, No.~5. --- P.~977--981.

\bibitem{Rosner1978}
Rosner, R., Tucker, W.~H., Vaiana, G.~S. Dynamics of the quiescent solar corona // \emph{The Astrophysical Journal}. --- 1978. --- Vol.~220. --- P.~643--665.

\bibitem{Klimchuk2008}
Klimchuk, J.~A., Patsourakos, S., Cargill, P.~J. Highly Efficient Modeling of Dynamic Coronal Loops // \emph{The Astrophysical Journal}. --- 2008. --- Vol.~682, No.~2. --- P.~1351--1362.

\bibitem{Lionello2009}
Lionello, R., Linker, J.~A., Miki\'c, Z. Multispectral Emission of the Sun During the First Whole Sun Month: Magnetohydrodynamic Simulations // \emph{The Astrophysical Journal}. --- 2009. --- Vol.~690, No.~1. --- P.~902--912.

\bibitem{HLL1983}
Harten, A., Lax, P.~D., van Leer, B. On upstream differencing and Godunov-type schemes for hyperbolic conservation laws // \emph{SIAM Review}. --- 1983. --- Vol.~25, No.~1. --- P.~35--61.

\end{thebibliography}

\end{document}